\section{问题一的建模与求解}

本节先分析问题一的输入、输出和关键难点，再根据题型完成数据处理、模型建立、求解和结果解释，最后给出直接回答题意的结论。

\subsection{问题分析}
说明问题一要回答的业务问题、数据粒度、关键难点和拟采用的建模路线。

\subsection{数据预处理与特征构造}
说明缺失值、异常值、极端值、单位和分析粒度处理，并给出关键指标或特征定义。

\subsection{模型建立与类别/参数选择}
\subsubsection{变量或指标定义}
说明变量、指标、标准化和方向。

\subsubsection{候选模型与评价标准}
说明候选模型、类别数或参数范围及选择依据。

\subsection{模型求解与结果分析}
\subsubsection{主模型结果}
呈现核心结果表、图和关键数值。

\subsubsection{业务画像或机制解释}
解释结果如何回答问题一，并给出可执行含义。

\subsection{问题一小结}
用简短结论回答问题一，保留关键分类、关联强度或排序结果，不重复算法细节。
